\documentclass[10pt,a4paper]{amsart}
\usepackage[margin=19mm]{geometry}
\usepackage{amsmath,amssymb,mathtools}
\usepackage[scr=rsfs,cal=euler]{mathalfa}
\usepackage{xfrac}
\usepackage[unicode,hidelinks,bookmarks=false]{hyperref}
\newcommand{\ie}{i.e.\ }
\newcommand{\cf}{cf.\ }
\newcommand{\resp}{respectively\ }
\newcommand{\Subsection}{\subsection}
\providecommand{\nobreakdash}{\nobreak\hbox{-}}
\setlength{\emergencystretch}{2em}
\multlinegap0pt
\usepackage{amssymb}
\usepackage{float}
\newtheorem{lem}{Lemma}
\newcommand{\beas}{\begin{eqnarray*}}
\newcommand{\eeas}{\end{eqnarray*}}
\title{Geometric analysis of a class of harmonic mappings defined by a differential Inequality}
\author{Vasudevarao Allu, Raju Biswas, Rajib Mandal}
\date{Source: arXiv:2501.01663v4 (CC BY 4.0)}
\begin{document}
\maketitle

\noindent\emph{Source and licence.} Geometric analysis of a class of harmonic mappings defined by a differential Inequality. Version arXiv:2501.01663v4 (CC BY 4.0), distributed by arXiv under the Creative Commons Attribution 4.0 International licence, http://creativecommons.org/licenses/by/4.0/, as recorded in the arXiv licence metadata for this exact version and shown on the abstract page https://arxiv.org/abs/2501.01663v4. Full licence notice: \texttt{licenses/arxiv-2501.01663-CC-BY-4.0.txt}.\par\smallskip

\noindent\textit{Editorial scope.} Counted unit: the article's lem lem1 (source0.tex lines 456--458). The article's complete original proof of it is delivered with it (source0.tex lines 459--471, one \texttt{proof} environment of 13 lines). No mathematics is written, repaired or re-derived: the statement and the proof are the article's own lines, in the article's own notation, and no retained line is substituted or re-flowed.  No further source-local context is needed: every cross-reference of every retained line resolves inside the two delivered spans.  Nothing was omitted from any retained span, so the package carries no omission record.  No retained line contains a citation, so the wrapper carries no bibliography and no bibliography entry.  No retained line contains a figure environment, an image-inclusion command or drawing code, so no image is delivered and the package declares no asset dependency.  Editorial formatting only, in addition: an AMS \texttt{amsart} preamble that restates the article's own declarations and macros used by the retained lines, namely the environments \texttt{lem} on separate counters, together with the standard packages those lines need.  The retained environments print in this document as Lemma 1 in order of appearance, whereas the article prints the same items with the numbering of its own full document.  The complete native source of the article as distributed in the arXiv e-print for this exact version is bundled unchanged as \texttt{sources/171102/source0.tex}.
\begin{lem}\label{lem1}
The class \(\mathcal{P}(\alpha, M)\) is closed under convex combinations.
\end{lem}

\begin{proof}
Let $F_1, F_2, \ldots, F_n \in \mathcal{P}(\alpha, M)$. Consider the convex combination $F=\sum_{i=1}^n t_i F_i$, where $t_i\geq 0$ for each $i$ with $\sum_{i=1}^n t_i=1$.
Thus, for each $i$, we have
\beas \text{Re}\left((1-\alpha)F_i'(z) + \alpha z F_i''(z)\right) > -M \quad \text{for }\quad z \in \mathbb{D}.\eeas
Therefore.
\beas
\text{Re}\left((1-\alpha)F'(z) + \alpha z F''(z)\right)
&= & \text{Re}\left(\sum_{i=1}^n t_i \left((1-\alpha)F_i'(z) + \alpha z F_i''(z)\right)\right) \\[1mm]
&=& \sum_{i=1}^n t_i  \text{Re}\left((1-\alpha)F_i'(z) + \alpha z F_i''(z)\right) \\[1mm]
&>& \sum_{i=1}^n t_i (-M) = -M,
\eeas
which shows that $F \in \mathcal{P}(\alpha, M)$. This completes the proof.
\end{proof}

\end{document}
